\subsection{Verifying the convolution obstruction}\label{p:sec:obstruction}

It remains to exclude invariant projections. We verify the hypotheses
of Lemma~\ref{lem:conditional-convolution} for the product couplings.
The infinite-sum estimate is already contained in that lemma; the task
here is to bound the finite coefficients on an event of fixed positive
probability and to attenuate every omitted index.

The preceding sections have constructed the factor $M$, its bounded
operator $T$, and all the couplings, with the scales and angles fixed
independently of any projection.  Let $p\in M$ be an orthogonal projection satisfying
$(1-p)Tp=0$, and suppose, toward a contradiction, that $0<p<1$.
Use the common invariant Borel conull set for this fixed $p$ and the
cocycles.  Equation~\eqref{p:eq:strict-diagonal} gives
\[
 0<r(x)=p(x)_{0,0}<1\qquad\text{for }\mu\text{-almost every }x.
\]
Every coupling used below has absolutely continuous endpoint
marginals, so its endpoints belong to this set almost surely.  All
pointwise operator and cocycle statements below are understood on
these sets.  Only countably many indices occur, so they may be
required simultaneously.

\paragraph{Convolution identities.}
Set
\[
 \xi(x)=p(x)e_0,\qquad \gamma(y)=(1-p(y))e_0.
\]
These measurable vectors have norms at most one, with
$\xi_0(x)=r(x)$ and $\gamma_0(y)=1-r(y)$.  Fix $j$.  By
Proposition~\ref{p:prop:biased-identity}, the vector
\[
 \zeta=\operatorname{diag}_{v\in\Z}\bigl(R_{x,y}(v)\bigr)\xi(x)
\]
belongs to $\operatorname{ran}p(y)$ for $\nu_j$-almost every
$(x,y)$.  The same range is invariant under the fiber $T(y)$ and
therefore under all its nonnegative powers.  By the cocycle identity,
the coordinate $n+v$ of $T(y)^n\zeta$ is
\[
 F_n(\sigma^v y)\,R_{x,y}(v)\,\xi_v(x)
 =\frac{F_{n+v}(y)}{F_v(x)}\xi_v(x),\qquad n\ge0.
\]
Taking its inner product with $\gamma(y)$ yields
\begin{equation}\label{p:eq:convolution-identity}
 0=\sum_{v\in\Z}\frac{F_{n+v}(y)}{F_v(x)}
                  \xi_v(x)\overline{\gamma_{n+v}(y)}
   =\sum_{v\in\Z}A_vB_{-n-v},\qquad n\ge0,
\end{equation}
where
\begin{equation}\label{p:eq:coefficient-sequences}
 A_v=\frac{\xi_v(x)}{F_v(x)},\qquad
 B_v=F_{-v}(y)\overline{\gamma_{-v}(y)}.
\end{equation}
These sequences themselves need not be square summable.  Every
convolution in \eqref{p:eq:convolution-identity}, however, is absolutely
convergent.  Indeed $F_{k+1}(z)=f(\sigma^k z)F_k(z)\le F_k(z)$
for every integer $k$, so for $n\ge0$,
\begin{equation}\label{p:eq:two-ratio-bounds}
 0<\frac{F_{n+v}(y)}{F_v(x)}
 \le\min\{R_{x,y}(v),R_{x,y}(n+v)\}\le1.
\end{equation}
The first bound uses $F_{n+v}(y)\le F_v(y)$ and the second uses
$F_v(x)\ge F_{n+v}(x)$.  Cauchy--Schwarz now bounds the absolute
sum by $\|\xi(x)\|_2\|\gamma(y)\|_2\le1$.

\paragraph{An event of uniformly positive probability.}
Write $K=K_j$ and $m=m_j$.  On the augmented probability space
$(\Omega_j,\mathbb P_j)$, define $E_j$ by the following conditions:
\begin{enumerate}
 \item For every $i\in J_j$,
 $x_i<1-2/m$ and $y_i>2/m$, and at least one coordinate in $J_j$
 has the tagged label.
 \item For every $s>j$ and $i\in J_s$,
 $x_i<1-2K/N_s$ and $y_i>2K/N_s$.
 \item $r(x)\ge1/j$ and $1-r(y)\ge1/j$.
\end{enumerate}
We claim that there is a constant $c_*>0$ such that
\begin{equation}\label{p:eq:positive-event}
 \mathbb P_j(E_j)\ge c_*
 \qquad\text{for all sufficiently large }j.
\end{equation}
The constant can be chosen independently of $p$; the point beyond
which the bound holds may depend on $p$.

At one coordinate of $J_j$, a tagged choice has probability $1/m$
and satisfies both of the first condition's inequalities.  By the
one-coordinate density bound in Lemma~\ref{p:lem:coupling-marginals},
the probability of a nontagged choice satisfying those inequalities
is at least
\[
 1-\frac1m-10\frac2m-10\frac2m=1-\frac{41}{m}.
\]
The disjoint events specifying exactly one tag and admissible
nontagged choices at the other coordinates consequently give the
lower bound
\begin{equation}\label{p:eq:tag-event}
 m\frac1m\left(1-\frac{41}{m}\right)^{m-1}
 =\left(1-\frac{41}{m}\right)^{m-1}\ge e^{-82}.
\end{equation}
Here $m\ge100$ and $\log(1-t)\ge-2t$ for $0\le t\le1/2$.
The probability that the second condition fails is at most
\begin{equation}\label{p:eq:future-event-loss}
 40K\sum_{s>j}\frac{m_s}{N_s}
 =40K\sum_{s>j}\frac1{K_s}
 \le\frac{80K_j}{K_{j+1}}\longrightarrow0.
\end{equation}
This is a countable union bound with a convergent sum.  The intervals
involved have length $2K/N_s$; all are less than one by the scale
conditions.

The two endpoint losses in the third condition are bounded, using
\eqref{p:eq:uniform-absolute-continuity}, by
\[
 \mathbb P_j\{r(x)<1/j\}\le H\mu\{r<1/j\}^{1/2},
 \qquad
 \mathbb P_j\{1-r(y)<1/j\}\le H\mu\{1-r<1/j\}^{1/2}.
\]
Both tend to zero because $0<r<1$ almost everywhere.  Subtracting
these losses and \eqref{p:eq:future-event-loss} from
\eqref{p:eq:tag-event} proves \eqref{p:eq:positive-event}, for example
with $c_*=e^{-82}/2$.  No independence between these conditions is
needed.

\paragraph{The finite coefficients.}
On $E_j$ the coefficients in \eqref{p:eq:coefficient-sequences} satisfy
\begin{equation}\label{p:eq:middle-coefficients}
 |A_v|,|B_v|\le D_j\quad(|v|\le K),
 \qquad |A_0|,|B_0|\ge1/j.
\end{equation}
To verify all signs, first let $0\le v\le K$.  For $i\in J_j$,
$v\theta_i\le2/m$, and for $i\in J_s$ with $s>j$,
$v\theta_i\le2K/N_s$.  The defining endpoint exclusions in $E_j$
therefore give
\[
 \lfloor x_i+v\theta_i\rfloor=0,
 \qquad \lfloor y_i-v\theta_i\rfloor=0
 \quad(i\in J_s,\ s\ge j).
\]
At each of the finitely many earlier coordinates, these floors have
absolute value at most $K+1$, since $0<\theta_i<1$.  Thus
\[
 \frac1{F_v(x)}\le
 \exp\left((K+1)\sum_{s<j}C_sm_s\right)=D_j,
 \qquad F_{-v}(y)\le D_j,
\]
which bounds $A_v$ and $B_v$.  For the negative coefficient indices,
monotonicity gives $F_{-v}(x)\ge1$ and $F_v(y)\le1$, and hence
\[
 |A_{-v}|\le1,\qquad |B_{-v}|\le1\qquad(0\le v\le K).
\]
Finally $A_0=r(x)$ and $B_0=1-r(y)$, proving
\eqref{p:eq:middle-coefficients}.

Thus the finite-coefficient hypotheses of
Lemma~\ref{lem:conditional-convolution} hold with $D=D_j$ and $a=1/j$.
We next verify its attenuation hypothesis outside this window.

\paragraph{Attenuation at every omitted index.}
We next control $R_{x,y}(k)$ when $|k|>K$.  If
$K<|k|\le N_j/4$, choose a tagged coordinate in $J_j$ on $E_j$.
At that coordinate
\[
 \frac1m<|k|\theta_i\le\frac12,
 \qquad x_i<\frac1m,\qquad y_i\ge1-\frac1m.
\]
For $k>0$ the threshold $1-\{k\theta_i\}$ lies in
$[1/2,1-1/m)$; for $k<0$ it equals $|k|\theta_i$ and lies in
$(1/m,1/2]$.  The tagged pair crosses it in both cases.  The
coordinate contribution to the ratio gives
\begin{equation}\label{p:eq:tag-attenuation}
 R_{x,y}(k)\le e^{-C_j}\le e^{-L_j}
 \quad\text{on }E_j,\qquad K<|k|\le N_j/4.
\end{equation}

Fix $j$ sufficiently large that the phase conclusions of
Proposition~\ref{p:prop:phases} hold in every block $s\ge j$.
For $|k|>N_j/4$ we can assign a block $J_s$ and a width $w$ so that
either the second phase range is used with $s\ge j$, $w=b_s$, or
the first phase range is used with $s>j$, $w=a_s$.  Here is the
coverage explicitly.  Put
\[
 \alpha_s=G_sK_s,\qquad \beta_s=N_s/4,
 \qquad Q_s=m_s^{(4/5)m_s}.
\]
The scale conditions imply
\[
 \alpha_s\le\beta_s<Q_s,\qquad \alpha_{s+1}<Q_s,
 \qquad Q_s\longrightarrow\infty.
\]
Start with $[\beta_j,Q_j]$, where the second phase range applies.
For each $s>j$, the first and second phase ranges together cover
$[\alpha_s,Q_s]$, divided at $\beta_s$.  Consecutive intervals
overlap by $\alpha_{s+1}<Q_s$, so their union with the initial
interval covers $[\beta_j,\infty)$.  These are real intervals;
in particular every integer magnitude in question is covered.

Choose one such assignment deterministically for each $k$, and let
$I(k)\subseteq J_s$ be a deterministic list of qualifying
coordinates.  The choice can, for example, use the least available
block and then increasing coordinate order.  With $q=|I(k)|$ we have
\begin{equation}\label{p:eq:tail-thresholds}
 q\ge m_s/20,\qquad
 u_i:=1-\{k\theta_i\}\in[w,1-w]\quad(i\in I(k)).
\end{equation}
All choices depend only on $j,k$ and the already fixed parameters and
angles, not on the sampled endpoints.  In the second phase range,
\[
 m_s\rho_s(b_s)=\kappa_sm_s^{1/40}\ge L_s\ge L_j;
\]
in the first phase range, where $s>j$,
\[
 m_s\rho_s(a_s)=\kappa_s\sqrt{G_s}\ge L_{s-1}\ge L_j.
\]
Also $C_s\ge L_s\ge L_j$ in either case.  If any listed threshold
is crossed, then
\begin{equation}\label{p:eq:tail-crossing-attenuation}
 R_{x,y}(k)\le e^{-C_s}\le e^{-L_j}.
\end{equation}

Let $\mathcal B_k$ be the event that none of the thresholds in
\eqref{p:eq:tail-thresholds} is crossed, and define
\[
 a_k(x)=\#\{i\in I(k):x_i<u_i\},
\]
\[
 \begin{aligned}
 H_k^{(1)}&=\mathcal B_k\cap\{a_k(x)\ge q/2\},\\
 H_k^{(2)}&=\mathcal B_k\cap\{a_k(x)<q/2\}.
 \end{aligned}
\]
Given $x$, the eligible coordinates in $H_k^{(1)}$ are fixed.
Conditional independence and \eqref{p:eq:crossing-given-x} imply
\begin{equation}\label{p:eq:first-bad-probability}
 \mathbb P_j(H_k^{(1)}\mid x)
 \le(1-\rho_s(w)/10)^{q/2}
 \le\exp(-m_s\rho_s(w)/400)
 \le e^{-L_j/400}.
\end{equation}
For the other orientation, orderedness gives on $\mathcal B_k$ the
exact equivalence
\[
 x_i<u_i\quad\Longleftrightarrow\quad y_i<u_i
 \qquad(i\in I(k)).
\]
Indeed the forward implication follows from the absence of a
crossing, and the reverse implication from $x_i\le y_i$.  Thus
\begin{equation}\label{p:eq:second-bad-event}
 H_k^{(2)}
 =\mathcal B_k\cap
   \left\{\#\{i\in I(k):y_i\ge u_i\}>q/2\right\}
 \quad\text{almost surely}.
\end{equation}
This is the same complementary conditioning split shown in
Figure~\ref{fig:conditioning}, with independent coordinates in place
of digit labels.
The eligible coordinates in this expression are determined by $y$.
The reverse conditional independence and
\eqref{p:eq:crossing-given-y} give
\begin{equation}\label{p:eq:second-bad-probability}
 \mathbb P_j(H_k^{(2)}\mid y)\le e^{-L_j/400}.
\end{equation}
These estimates hold under the full probability law $\mathbb P_j$,
before any restriction to $E_j$.  There are countably many $k$, so
the conditional estimates can be taken simultaneously outside
endpoint null sets.  No union estimate over their bad events has
been made.


For the near-window indices in \eqref{p:eq:tag-attenuation}, set both
exceptional events equal to the empty set. Combining that estimate with
\eqref{p:eq:tail-crossing-attenuation} and the two conditional bounds, we
have, for every $|k|>K_j$,
\[
 R_{x,y}(k)\le e^{-L_j}
       +\ind_{H_k^{(1)}}+\ind_{H_k^{(2)}}\quad\text{on }E_j,
\]
where
\[
 \mathbb P_j(H_k^{(1)}\mid x)\le e^{-L_j/400},\qquad
 \mathbb P_j(H_k^{(2)}\mid y)\le e^{-L_j/400}
\]
hold under the full law. All inputs to
Lemma~\ref{lem:conditional-convolution} are now verified: the vectors,
convolutions and ratio bounds are those of
\eqref{p:eq:convolution-identity}--\eqref{p:eq:two-ratio-bounds}; the event
and finite coefficient bounds are \eqref{p:eq:positive-event} and
\eqref{p:eq:middle-coefficients}. For every sufficiently large fixed $j$,
the lemma therefore gives
\begin{equation}\label{p:eq:final-contradiction}
 \begin{split}
 c_*\varepsilon_j
 &\le (2K_j+1)\bigl(e^{-L_j}+2e^{-L_j/800}\bigr)\\
 &\le (2K_j+1)3e^{-L_j/10000}
 \le \frac{\varepsilon_j}{j},
 \end{split}
\end{equation}
by \eqref{p:eq:epsilon-choice}. Since $\varepsilon_j>0$, cancellation at
this fixed stage gives $c_*\le1/j$, impossible once $j>1/c_*$.
No uniform lower bound on $(\varepsilon_j)$ is used. The scales, angles,
operator and coupling laws were fixed before the arbitrary projection
$p$; only its conull representatives and the eventual threshold in
\eqref{p:eq:positive-event} depended on $p$. This excludes every
nontrivial invariant projection in $M$.
